\documentclass[11pt,a4paper]{article}
\usepackage[margin=27mm]{geometry}
\usepackage[T1]{fontenc}\usepackage{lmodern}
\usepackage{amsmath,amssymb,amsthm,mathtools,booktabs,microtype}
\usepackage[colorlinks=true,linkcolor=blue,citecolor=blue,urlcolor=blue]{hyperref}
\newtheorem{theorem}{Theorem}
\newtheorem{lemma}[theorem]{Lemma}
\newtheorem{corollary}[theorem]{Corollary}
\theoremstyle{remark}\newtheorem{remark}[theorem]{Remark}
\DeclareMathOperator{\rank}{rank}\DeclareMathOperator{\Tr}{Tr}
\newcommand{\HH}{\mathcal H}\newcommand{\K}{\mathcal K}
\newcommand{\DD}{\mathcal D}\newcommand{\norm}[1]{\left\|#1\right\|_{\!F}}
\title{A sharp stability transition for low-rank tomography\\with four missing Weyl coefficients}
\author{Lluis Eriksson\\\small Independent researcher, Sweden}
\date{24 October 2026\\Working draft 2; documentary clarifications}
\begin{document}\maketitle
\begin{abstract}
Let $d=9n$ with arbitrary integer $n\ge1$, and omit the four cyclic Weyl
coefficients associated with $X_d^{d/3}$, its adjoint, $Z_d^{d/3}$ and its
adjoint. We determine the exact inverse Frobenius constant on differences
of density matrices of rank at most $r$. Writing $x=r/n$, the squared
fraction of energy that can be hidden in the missing coefficients is
$2x/3$ for $0<x\le1$, and
$(x+\sqrt{x^2-3x+3})/3$ for $1\le x<2$.
The resulting inverse constant is attained by explicit orthogonally
supported states and diverges as $\sqrt{2n/(2n-r)}$ at the exact
nonidentifiability threshold $r=2n$. The same threshold holds when a
rank-$r$ state must be distinguished from every positive state of
unrestricted rank. For this stronger task we give a separate, conservative
stability bound and a convex least-squares recovery consequence. The
proof reduces the missing-coefficient space to additive $3$ by $3$
spectral grids, then solves two four-dimensional symmetric matrix
pencils. These are results for one fixed erasure pattern, not for
arbitrary four erasures or a measurement scheme with subquadratically
many coefficients.
\end{abstract}

\section{Question, normalization and antecedents}
A measurement may fail to determine arbitrary quantum states while still
determining every state below a rank threshold. In that situation a
quantitative question remains: how much can the omitted data amplify
errors even before injectivity fails? We solve that question exactly for
a four-dimensional missing-coefficient space.

Put $\omega_d=e^{2\pi i/d}$ and
\[
 X_de_j=e_{j+1\bmod d},\qquad Z_de_j=\omega_d^je_j,\qquad
 W_d(a,b)=X_d^aZ_d^b .
\]
The $d^2$ Weyl matrices are orthogonal, with squared Frobenius norm $d$.
Throughout, $d=9n$ and $n$ is a positive integer. Define
\[
 U=X_d^{d/3},\qquad V=Z_d^{d/3},\qquad
 E=\{U,U^\dagger,V,V^\dagger\}.
\]
These are four distinct nonidentity Weyl matrices. Let $P$ be the
Frobenius-orthogonal projection onto their complex span, and put $Q=I-P$.
On Hermitian matrices $P$ is the projection onto the real space
\begin{equation}\label{eq:kernel}
 \K=\{aU+\bar a U^\dagger+bV+\bar b V^\dagger:a,b\in\mathbb C\}.
\end{equation}
It consists of traceless Hermitian matrices. Thus $Q\rho$ means all Weyl
coefficients except these four, including the known trace coefficient.
For $C=\sum_w c_wW_w$ our data norm is
$\norm{QC}^{\,2}=d\sum_{w\notin E}|c_w|^2$.
All constants below refer to this normalization, not directly to a
statistical probability norm.

This coefficient space can be realized by a POVM, although outcome norms
then require a change of coordinates. Choose a real Hermitian basis
$F_1,\ldots,F_m$ of the traceless part of $\K^\perp$, where $m=d^2-5$,
and rescale each element so that $\|F_j\|_{\rm op}\le1$.
The $2m$ effects $(I+F_j)/(2m),(I-F_j)/(2m)$ are positive, sum to $I$,
and span $\K^\perp$. Their probability record determines $Q\rho$.
The constants below do not include the conditioning of this coordinate
conversion or the sampling error of this particular realization.

Let $\DD_r=\{\rho\ge0:\Tr\rho=1,\rank\rho\le r\}$ and define
\begin{equation}\label{eq:C}
 C(n,r)=\sup_{\substack{\rho,\sigma\in\DD_r\\\rho\ne\sigma}}
       \frac{\norm{\rho-\sigma}}{\norm{Q(\rho-\sigma)}} ,
\end{equation}
where a zero denominator gives infinity. Ranks are integers.

The general kernel and inertia characterization of restricted quantum
state determination is due to Carmeli, Heinosaari, Schultz and
Toigo~\cite{carmeli}. Their Section 6.2, equations (12)--(13) and
Proposition 13, also establish the relation between symmetric Weyl zero
sets and realizable covariant POVM kernels. The elementary realization
above is not claimed as a new measurement-design principle.
Baldwin, Deutsch and Kalev~\cite{bdk}, especially
Sections II--III, distinguish rank-$r$ completeness from
rank-$r$ strict-completeness and explain the latter's relation to convex
reconstruction and noise robustness. We use this established distinction.
We supply direct proofs of the specialized inertia and stability claims
rather than assert a new general criterion.

Recent work of Zhu and Wang~\cite{zhu} studies spectral floors and
conditioning for minimal rank-one Weyl--Heisenberg covariant
measurements. Their full-space Gram-spectrum problem differs from the
restricted state-difference constant of the singular coefficient
projection considered here. The author's earlier Weyl rank-uncertainty
reductions~\cite{eriksson} concern support versus arbitrary complex
matrix rank, rather than the exact positive-state conditioning law
proved below. Spectral rearrangement, orthogonal projection duality and
symmetric matrix pencils are classical tools; the contribution claimed
here is their explicit solution for the kernel~\eqref{eq:kernel}.
This comparison is not a certification of worldwide priority.

\section{Exact stability and identifiability}
\begin{theorem}[Sharp balanced-rank stability]\label{thm:sharp}
For $1\le r<2n$ let $x=r/n$ and
\[
 L(x)=
 \begin{cases}
  2x/3,&0<x\le1,\\[2pt]
  \bigl(x+\sqrt{x^2-3x+3}\bigr)/3,&1\le x<2 .
 \end{cases}
\]
Then
\begin{equation}\label{eq:main}
 C(n,r)=(1-L(r/n))^{-1/2}.
\end{equation}
The supremum is attained by differences of states of rank exactly $r$
with mutually orthogonal supports. For $r\ge2n$ it is infinite.
In particular $C(n,n)=\sqrt3$, and along integer sequences with
$(2n-r)/n\to0$ from above,
\[
 C(n,r)\sim\sqrt{\frac{2n}{2n-r}}.
\]
\end{theorem}

\begin{theorem}[Unrestricted positive competitors]\label{thm:strict}
Every nonzero $H\in\K$ has at least $2n$ strictly positive and at least
$2n$ strictly negative eigenvalues. Consequently $Q$ distinguishes every
state of rank at most $r$ from every other density matrix, of any rank,
if and only if $r<2n$. In this range, for $\rho\in\DD_r$ and any density
matrix $\sigma$,
\begin{equation}\label{eq:strict}
 \norm{\rho-\sigma}\le B(n,r)\norm{Q(\rho-\sigma)},\qquad
 B(n,r)=\sqrt{1+\frac{288n}{2n-r}} .
\end{equation}
The bound $B$ is conservative; it is not asserted to be the optimal
constant for this larger competitor class.
\end{theorem}

The finite constants $C$ and $B$ have different domains. In particular,
\eqref{eq:main} cannot be substituted into \eqref{eq:strict} merely
because strict-completeness holds. For $n=1$ the finite range contains
only $r=1$; the intermediate branch becomes nontrivial for $n\ge2$.

\section{Spectral grids and the state-difference cone}
The unitaries $U,V$ commute, have order three, and generate nine distinct
Weyl labels. Every nonidentity $U^iV^j$, $0\le i,j<3$, has trace zero.
Character orthogonality therefore shows that each of the nine joint
eigenspaces has dimension $n$.

By the real three-point Fourier transform, every $H\in\K$ has joint
eigenvalues
\begin{equation}\label{eq:grid}
 h_{ij}=f_i+g_j,\qquad
 \sum_{i=1}^3 f_i=\sum_{j=1}^3g_j=0,
\end{equation}
each repeated $n$ times, where $f,g$ are arbitrary real zero-mean
triples. Conversely every such pair occurs in \eqref{eq:kernel}. In
particular
\begin{equation}\label{eq:gridnorm}
 \norm H^2=3n\left(\sum_i f_i^2+\sum_jg_j^2\right).
\end{equation}
No restriction to dimensions that are powers of three is needed.

Write $n_+(D),n_-(D)$ for the positive and negative inertias of a
Hermitian matrix. Up to a positive scalar, the nonzero differences of
states in $\DD_r$ are exactly
\[
 \mathcal C_r=\{D=D^\dagger\ne0:\Tr D=0,\ n_+(D),n_-(D)\le r\}.
\]
Indeed, a state difference has these inertia bounds by the min--max
principle. Conversely the positive and negative spectral parts of such
a $D$ have a common nonzero trace $\tau$; $D/\tau$ is the difference of
their trace-one normalizations. Hence the key extremum is
\begin{equation}\label{eq:energy}
 L_*=\sup_{D\in\mathcal C_r}\frac{\norm{PD}^2}{\norm D^2}.
\end{equation}
Orthogonality gives $C(n,r)=(1-L_*)^{-1/2}$ whenever $L_*<1$.

\begin{lemma}[A common-shift dual bound]\label{lem:dual}
For an integer $1\le r<2n$ and $H=H^\dagger$ with eigenvalues
$h_1\ge\cdots\ge h_d$, put
\[
 W_r(H)=\min_{\theta\in\mathbb R}
       \sum_{i=1}^r\bigl[(h_i-\theta)^2+(h_{d-i+1}-\theta)^2\bigr].
\]
For $D\in\mathcal C_r$, $|\Tr(HD)|^2\le W_r(H)\norm D^2$. Thus
\begin{equation}\label{eq:dual}
 L_*\le\sup_{0\ne H\in\K}\frac{W_r(H)}{\norm H^2}.
\end{equation}
\end{lemma}
\begin{proof}
Pad the positive eigenvalues $p_i$ of $D$ and negative magnitudes $q_i$
with zeros to length $r$, ordering both decreasingly. Their sums agree.
Trace rearrangement bounds $\Tr(HD)$ by
$\sum_{i\le r}h_i p_i-\sum_{i\le r}h_{d-i+1}q_i$.
For completeness, squared overlaps of two orthonormal eigenbases give a
doubly stochastic matrix; the extremal pairing of two ordered real
lists bounds the trace. The same conclusion follows from weighted
spectral variational inequalities. This argument does not require $D$
to commute with $H$.

Subtract $\theta(\sum p_i-\sum q_i)=0$ and apply Cauchy--Schwarz.
Apply the same reasoning to $-H$ to control the absolute value, noting
$W_r(-H)=W_r(H)$. Finally
$\norm{PD}=\max_{H\in\K,\norm H=1}|\Tr(HD)|$, which proves
\eqref{eq:dual}.
\end{proof}

\section{Solving the dual extremum}
\subsection{Ranks at most one joint multiplicity}
For $r\le n$ the extreme values of $H$ each occur at least $n$ times, so
\[
 W_r(H)=\frac r2(\max h-\min h)^2.
\]
For a zero-mean triple,
$\operatorname{diam}(f)^2\le2\sum f_i^2$. Since
$\operatorname{diam}(h)=\operatorname{diam}(f)+\operatorname{diam}(g)$,
\eqref{eq:gridnorm} gives $W_r(H)/\norm H^2\le2r/(3n)$.

Take $f=(1,0,-1)$ and $g=-f$. The grid has extrema $\pm2$ on spaces
of dimension $n$, and $\norm H^2=12n$. Choose rank-$r$ projectors
$R_+,R_-$ in those two spaces and set $D=2(R_+-R_-)$.
Then $\Tr(HD)=\norm D^2=8r$. Projection duality yields
\[
 \frac{\norm{PD}^2}{\norm D^2}\ge
 \frac{|\Tr(HD)|^2}{\norm H^2\norm D^2}
 =\frac{2r}{3n}.
\]
This matches the upper bound, and $R_+/r,R_-/r$ are the required states.

\subsection{Ranks between one and two joint multiplicities}
Now $n\le r<2n$, $x=r/n$, and $z=x-1$. Sort $f,g$ decreasingly and put
\[
 a=f_1-f_2,\quad b=f_2-f_3,\quad c=g_1-g_2,\quad e=g_2-g_3.
\]
These gaps are nonnegative. Define
\[
 K=a^2+ab+b^2+c^2+ce+e^2,\quad
 A=a+b+c+e,\quad u=\min(a,c),\quad v=\min(b,e).
\]
Then $\norm H^2=2nK$. The largest two cells are
$h_{\max},h_{\max}-u$ and the smallest two
$h_{\min},h_{\min}+v$, with weights $n,r-n$ at each end.
Ties and zero gaps cause no problem: consistent cell ordering selects
the same multiset, with partial cells of dimension $r-n$.

The shift minimizing $W_r$ is
\[
 \theta=\frac{h_{\max}+h_{\min}}2+\frac{z(v-u)}{2x}.
\]
Expanding the squares gives
\begin{equation}\label{eq:Q}
 \frac{W_r(H)}n=
 \frac{xA^2}2-zA(u+v)+z(u^2+v^2)-\frac{z^2(v-u)^2}{2x}.
\end{equation}
The single common shift is essential because the positive and negative
parts of a state difference must have the same trace.

There are four choices for the two minima; exchange of $f,g$ reduces
them to $(u,v)=(a,b)$ and $(a,e)$. Write $w=(a,b,c,e)^T$ and
\[
 M=\begin{pmatrix}2&1&0&0\\1&2&0&0\\0&0&2&1\\0&0&1&2\end{pmatrix}>0.
\]
Then $2K=w^TMw$, and the numerator in \eqref{eq:Q} is $w^TNw$, with
\begin{align*}
 N_{\rm same}&=
 \begin{pmatrix}
 1-\frac1{2x}&\frac1{2x}&\frac12&\frac12\\
 \frac1{2x}&1-\frac1{2x}&\frac12&\frac12\\
 \frac12&\frac12&\frac x2&\frac x2\\
 \frac12&\frac12&\frac x2&\frac x2
 \end{pmatrix},\\
 N_{\rm cross}&=
 \begin{pmatrix}
 1-\frac1{2x}&\frac12&\frac12&\frac1{2x}\\
 \frac12&\frac x2&\frac x2&\frac12\\
 \frac12&\frac x2&\frac x2&\frac12\\
 \frac1{2x}&\frac12&\frac12&1-\frac1{2x}
 \end{pmatrix}.
\end{align*}
Their exact pencil determinants are
\begin{align}
 \det(N_{\rm same}-\lambda M)
 &=-\frac{\lambda}{x}(\lambda x-x+1)
       [-9\lambda^2+3(x+1)\lambda-x+1],\label{eq:det1}\\
 \det(N_{\rm cross}-\lambda M)
 &=-\frac{\lambda}{x}(3\lambda x-2x+2)
       [-3\lambda^2+2x\lambda-x+1].\label{eq:det2}
\end{align}
These identities can be checked by direct four-by-four expansion.
As $M$ is positive definite, the largest generalized eigenvalue bounds
the quotient even after relaxing the gap and branch inequalities.

Put $\ell=(x+\sqrt{x^2-3x+3})/3$. It is the larger root of
$3\lambda^2-2x\lambda+x-1=0$ and is at least $2/3$ for $1\le x\le2$.
The other roots of \eqref{eq:det2} are $0$,
$2(x-1)/(3x)\le1/3$, and the smaller quadratic root. Thus the cross
branch has maximum generalized eigenvalue $\ell$.

For the same branch the additional roots are $0$ and $(x-1)/x\le1/2$.
With $q(\lambda)=9\lambda^2-3(x+1)\lambda+x-1$,
\[
 q(\ell)=(x-1)(3\ell-2)\ge0.
\]
Moreover $\ell$ lies to the right of the vertex $(x+1)/6\le1/2$.
It follows that the larger root of $q$ is at most $\ell$. This proves
$L_*\le\ell$ in all branches.

\section{Extremizing states and the singular boundary}
For $n\le r<2n$ put
\[
 t=1-x+\sqrt{x^2-3x+3},\qquad
 f=\left(\frac{2+t}3,\frac{t-1}3,-\frac{1+2t}3\right),\quad g=-f .
\]
Here $0<t\le1$, and
\begin{equation}\label{eq:t}
 t^2+2(x-1)t+x-2=0.
\end{equation}
The grid of $H$ has eigenvalues $0$ on $3n$ dimensions and
$\pm(1+t),\pm1,\pm t$, each listed sign repeated $n$ times.
At $t=1$ distinct cells may share an eigenvalue; the indicated cell
decomposition is still valid.

Retain the full $n$-dimensional cells of $\pm(1+t)$, rank-$(r-n)$
subspaces in the cells of $\pm1$, and zero elsewhere, obtaining $D$.
Then $\Tr D=0$, $n_\pm(D)=r$, and
\[
 \Tr(HD)=\norm D^2=2[n(1+t)^2+(r-n)],\qquad
 \norm H^2=2n[(1+t)^2+1+t^2].
\]
Using \eqref{eq:t},
\[
 \frac{\norm D^2}{\norm H^2}
 =\frac{(1+t)^2+x-1}{2(1+t+t^2)}
 =\frac{x+\sqrt{x^2-3x+3}}3=\ell.
\]
Projection duality forces $\norm{PD}^2/\norm D^2\ge\ell$, so equality
holds by the previously proved upper bound. The positive and negative
parts of $D$, divided by their common trace $n(1+t)+(r-n)$, are
explicit orthogonally supported rank-$r$ states.

At $r=2n$, the matrix
\[
 H_0=U+U^\dagger-V-V^\dagger
\]
has eigenvalues $3,-3$ each on $2n$ dimensions and zero elsewhere.
Its positive and negative spectral parts, each divided by $6n$, are
distinct rank-$2n$ states whose difference lies in $\K$. This proves
failure for all $r\ge2n$. For $x<2$, $L(x)<1$, completing
\eqref{eq:main}. Finally $L'(2^-)=1/2$, whence
$1-L(x)\sim(2-x)/2$ and the stated asymptotic follows.

\section{Strict-completeness and convex reconstruction}
\begin{proof}[Proof of Theorem~\ref{thm:strict}]
For a nonzero grid \eqref{eq:grid}, sort both triples and put
$F=f_1\ge0$, $G=g_1\ge0$, $S=F+G>0$. Zero mean and ordering give
$f_2\ge-F/2$, $g_2\ge-G/2$. The second largest cell satisfies
\[
 h_2=\max(F+g_2,f_2+G)
 \ge\max(F-G/2,G-F/2)\ge S/4>0.
\]
Thus at least two cells, each of dimension $n$, are strictly positive.
Apply the argument to $-H$ for the negative cells.

We also have $\sum f_i^2\le6F^2$: indeed
$f_2\in[-F/2,F]$ and the convex quadratic
$F^2+f_2^2+(F+f_2)^2$ is at most $6F^2$ on that interval.
Similarly $\sum g_j^2\le6G^2$. Consequently
\[
 \norm H^2\le18n(F^2+G^2)\le18nS^2.
\]
The largest $2n$ positive eigenvalues are therefore all at least
$\norm H/\sqrt{288n}$, and the same bound holds for the $2n$ largest
negative magnitudes.

Suppose $D=D^\dagger$, $\Tr D=0$, and either $n_+(D)\le r$ or
$n_-(D)\le r$. Set $H=PD$. If $H=0$ there is nothing to prove.
In the first case the Hermitian eigenvalue distance inequality bounds
$\norm{D-H}^2$ below by the sum of squares of the positive eigenvalues
of $H$ after its $r$ largest ones. To see this, sort the eigenvalues of
both matrices; $D$ has at most $r$ positive entries, so each subsequent
positive eigenvalue of $H$ contributes at least its square to their
squared Euclidean distance, which is bounded above by
$\norm{D-H}^2$. This last comparison follows from trace rearrangement:
if $d_i,h_i$ are the respective decreasing eigenvalue lists, then
\[
 \norm{D-H}^2=\norm D^2+\norm H^2-2\Tr(DH)
 \ge\sum_i(d_i-h_i)^2.
\]
Hence
\[
 \norm{QD}^2\ge\frac{2n-r}{288n}\norm{PD}^2.
\]
Apply the same argument to $-D,-H$ in the second case. Pythagoras now
gives $\norm D\le B(n,r)\norm{QD}$.

If $\rho\in\DD_r$ and $\sigma$ is any density matrix, the quadratic form
of $D=\rho-\sigma$ is nonpositive on $\ker\rho$. By min--max,
$n_+(D)\le r$, so \eqref{eq:strict} follows. It implies uniqueness for
$r<2n$. The states constructed from $H_0$ prove failure at and above
$2n$.
\end{proof}

\begin{corollary}[A convex noisy-data estimator]\label{cor:noise}
Let $1\le r<2n$, $\rho\in\DD_r$, and $y=Q\rho+e$ with
$\norm e\le\epsilon$. Every minimizer
\[
 \widehat\rho\in\arg\min_{\sigma\ge0,\ \Tr\sigma=1}\norm{Q\sigma-y}
\]
satisfies $\norm{\widehat\rho-\rho}\le2B(n,r)\epsilon$.
\end{corollary}
\begin{proof}
The density-matrix domain is nonempty, compact and convex, so a minimizer
exists. Its residual is at most that of $\rho$, hence at most $\epsilon$.
The triangle inequality gives
$\norm{Q(\widehat\rho-\rho)}\le2\epsilon$, and
Theorem~\ref{thm:strict} applies without a rank restriction on
$\widehat\rho$.
\end{proof}

This is a deterministic guarantee in coefficient space. It supplies no
experimental sampling bound, approximate-rank guarantee or running-time
analysis. In particular it does not assert an improvement over compressed
sensing protocols using many fewer measured coefficients.

\section{Scope, checks and remaining questions}
The argument depends only on the additive spectral kernel with nine
equal-multiplicity cells. It therefore applies after arbitrary unitary
conjugation and to any commuting order-three pair with all nine joint
multiplicities equal to $n$. The cyclic Weyl construction realizes this
kernel in every dimension divisible by nine.

Exact symbolic checks reproduce both pencil determinants and the
attainer identities, including a direct computation of projection energy.
Finite exact tests additionally compare sorted spectra with the formulas
in $2{,}400$ formula checks on $2{,}000$ distinct instances, using
$n=1,\ldots,5$ and $80$ nonzero gap vectors. The $400$ instances at $r=n$
are checked in both branches.
These tests supplement the universal proofs; they do not establish
unexamined dimensions by extrapolation.

The main open quantitative question here is the optimal constant in
\eqref{eq:strict}. Further questions concern unequal joint
multiplicities, other erasure patterns, and error bounds for approximately
low-rank states. The present result does not classify those problems.
Only four coefficients are removed out of $d^2$; the theorem isolates an
exact conditioning mechanism rather than a large measurement reduction.

\paragraph{AI and responsibility disclosure.}
The named author directed substantial OpenAI Codex assistance in
derivation, writing and exact computational checks. GPT-6 Astra xhigh
mathematical audits used shared filesystem and declared reused context.
No runtime isolation, independent human review, Lean formalization,
certified worldwide originality or predicted assessment score is claimed.
A full assessment of draft 1 is preserved separately. This draft makes
documentary clarifications; no new full assessment score is assigned to it.

\begin{thebibliography}{9}
\bibitem{carmeli}
C.~Carmeli, T.~Heinosaari, J.~Schultz and A.~Toigo,
\emph{Tasks and premises in quantum state determination},
arXiv:1308.5502v2 (2014).
\url{https://arxiv.org/abs/1308.5502v2}.
\bibitem{bdk}
C.~H.~Baldwin, I.~H.~Deutsch and A.~Kalev,
\emph{Strictly-complete measurements for bounded-rank quantum-state
tomography}, Physical Review A \textbf{93}, 052105 (2016).
\url{https://doi.org/10.1103/PhysRevA.93.052105};
arXiv:1605.02109v1.
\bibitem{zhu}
X.~Zhu and Y.~Wang,
\emph{Uniformly Stable Minimal Weyl--Heisenberg Measurements Approaching
the SIC Benchmark}, arXiv:2608.11850v1 (2026).
\url{https://arxiv.org/abs/2608.11850v1}.
\bibitem{eriksson}
L.~Eriksson, \emph{Removing the binary and ternary factors from cyclic
Weyl rank uncertainty}, AIRR.SCIENCE, ARR-2026-1QEPAQ29TN8FZRK5,
v1 (2026).
\url{https://airr.science/papers/ARR-2026-1QEPAQ29TN8FZRK5/versions/v1/}.
\end{thebibliography}
\end{document}
